\section{Erd\H{o}s Problem 624 --- ROUND-2}

\subsection*{1) ROUND-2 OBJECTIVE}

\textbf{Path pursued: (C) obstruction / gap-isolation.}

Round-1 gave only lower-bound information:
the counting lower bound $H(n)\ge \log_2 n$, Alon's argument that $H(2^k)\ge k+1$,
and the observation that fixed-size collision arguments leave too much slack when
$h=\log_2 n + C$.
In this round I do two things:

\begin{itemize}
\item I prove a \emph{barrier theorem} showing that a broad class of first-moment / averaging arguments
can never beat the trivial bound $n\le 2^h$.
\item I prove a new unconditional \emph{upper bound}
\[
H(n)\le \log_2 n + 2\log_2\log_2 n + O(1),
\]
which is strictly stronger than the Round-1 quoted $(3+o(1))\log_2\log_2 n$ upper bound.
\end{itemize}

So Round-2 both sharpens the obstruction on the lower-bound side and advances the upper-bound side.

\subsection*{2) Round-1 FOUNDATION USED}

I use the following Round-1 results as fixed background:

\begin{itemize}
\item the restatement of $H(n)$;
\item the basic counting lower bound $H(n)\ge \lceil\log_2 n\rceil$;
\item Alon's argument that $H(2^k)\ge k+1$;
\item the Round-1 diagnosis that the naive ``many collisions in one fibre'' idea does not by itself force
an unbounded additive gap.
\end{itemize}

\subsection*{3) NEW INSIGHT / TOOL (ROUND-2)}

There are two new tools.

\begin{itemize}
\item \textbf{Boolean-lattice colouring viewpoint.}
A witness $f:\mathcal P(X)\to X$ is exactly an $n$-colouring of the Boolean lattice $2^X$ such that
every principal ideal $2^Y$ with $|Y|=h$ is rainbow.

\item \textbf{Random-colouring upper bound plus an LP/first-moment barrier.}
The random-colouring argument yields the upper bound
$H(n)\le \log_2 n + 2\log_2\log_2 n + O(1)$.
Separately, any lower-bound argument obtained by taking nonnegative linear combinations of the basic
surjectivity inequalities for $h$-sets collapses to the trivial counting bound $n\le 2^h$.
\end{itemize}

\subsection*{4) ATTACK PLAN (ROUND-2)}

The exact gap after Round-1 was:

\begin{quote}
How can one go beyond the trivial bound $n\le 2^h$ and the constant-gap argument $H(2^k)\ge k+1$?
\end{quote}

Round-2 attacks this in two directions.

\begin{enumerate}
\item Prove that the most obvious first-moment averaging strategy is \emph{provably incapable}
of giving anything stronger than $n\le 2^h$.
\item Compensate by strengthening the known upper bound through a direct random construction.
\end{enumerate}

This does not solve the original problem, but it cleanly separates what still needs to happen:
any eventual proof of $H(n)-\log_2 n\to\infty$ must use structure beyond these first-moment inequalities.

\subsection*{5) WORK (ROUND-2)}

\paragraph{Reformulation.}
Let $X$ be an $n$-element set, and let $h$ be such that there exists
$f:\mathcal P(X)\to X$ with the property that every $Y\subseteq X$ of size $h$ satisfies
\[
\{f(A):A\subseteq Y\}=X.
\]
For each $x\in X$, define
\[
\mathcal F_x:=\{A\subseteq X:f(A)=x\}.
\]
Then $(\mathcal F_x)_{x\in X}$ is a partition of $\mathcal P(X)$, and the defining property of $H(n)$ is:

\begin{quote}
for every $x\in X$ and every $Y\in \binom{X}{h}$, there exists $A\in \mathcal F_x$ with $A\subseteq Y$.
\end{quote}

Equivalently, each colour class $\mathcal F_x$ hits every principal ideal $2^Y$ with $|Y|=h$.

\paragraph{Barrier theorem (first-moment / averaging methods).}
Fix $h$.
For each $x\in X$ and $Y\in \binom{X}{h}$ we have the basic surjectivity inequality
\[
\sum_{A\subseteq Y}\mathbf 1_{\{f(A)=x\}}\ge 1.
\tag{1}
\]
Take any family of nonnegative coefficients $(c_Y)_{Y\in \binom{X}{h}}$ that do \emph{not} depend on $x$.
Multiply (1) by $c_Y$ and sum over all $Y\in \binom{X}{h}$.
For a fixed $x$, this gives
\[
\sum_{A\subseteq X}\left(\sum_{\substack{Y\in \binom{X}{h}\\ A\subseteq Y}} c_Y\right)\mathbf 1_{\{f(A)=x\}}
\ \ge\ 
\sum_{Y\in \binom{X}{h}} c_Y.
\tag{2}
\]
Now sum (2) over all $x\in X$.
Because the sets $\mathcal F_x$ partition $\mathcal P(X)$, each subset $A\subseteq X$ contributes exactly once.
Hence
\[
n\sum_{Y\in \binom{X}{h}} c_Y
\le
\sum_{Y\in \binom{X}{h}} c_Y\sum_{A\subseteq Y}1
=
2^h\sum_{Y\in \binom{X}{h}} c_Y.
\]
If not all $c_Y$ vanish, we may divide by the positive scalar $\sum_Y c_Y$ and obtain
\[
n\le 2^h.
\]

Therefore:

\begin{quote}
\emph{Any lower-bound argument that is obtained solely by taking nonnegative linear combinations
of the basic inequalities (1), with coefficients independent of the colour $x$, can recover at most
the trivial counting bound $n\le 2^h$.}
\end{quote}

This rigorously rules out an entire first-moment class of approaches.

\paragraph{New upper bound.}
I now prove
\[
H(n)\le \log_2 n + 2\log_2\log_2 n + O(1).
\]

\paragraph{Theorem.}
For all sufficiently large $n$,
\[
H(n)\le \left\lceil \log_2 n + 2\log_2\log_2 n + 2\right\rceil.
\]

\paragraph{Proof.}
Let $X$ be an $n$-element set, and put
\[
L:=\log_2 n,\qquad
h:=\left\lceil L+2\log_2 L+2\right\rceil.
\]
Choose a random map $f:\mathcal P(X)\to X$ by assigning to each subset $A\subseteq X$
an independent uniformly distributed value $f(A)\in X$.

Fix $x\in X$ and $Y\in \binom{X}{h}$.
The event that $x$ is \emph{missing} from the image $\{f(A):A\subseteq Y\}$ is exactly the event that
all $2^h$ subsets of $Y$ avoid the value $x$, and therefore has probability
\[
\Pr\bigl(x\notin \{f(A):A\subseteq Y\}\bigr)
=
\left(1-\frac1n\right)^{2^h}
\le e^{-2^h/n}.
\]
By the union bound, the probability that \emph{some} pair $(x,Y)$ is bad is at most
\[
n\binom{n}{h}e^{-2^h/n}.
\tag{3}
\]
So it is enough to show that the quantity in (3) is $<1$.

Now $h\le 2L$ for all sufficiently large $n$.
Using the standard bound $\binom{n}{h}\le (en/h)^h$, we get
\[
\ln\!\left(n\binom{n}{h}\right)
\le
\ln n + h\ln\!\left(\frac{en}{h}\right)
\le
\ln n + 2L(\ln n+1)
=
2(\ln 2)L^2 + O(L).
\tag{4}
\]
Here the last inequality uses only $h\le 2L$ and the elementary bound
$\ln(en/h)\le \ln n+1$.
On the other hand, by the definition of $h$,
\[
\frac{2^h}{n}\ge 4L^2.
\tag{5}
\]
Since $4>2\ln 2$, the right-hand side of (5) is larger than the right-hand side of (4) for all sufficiently large $n$.
Hence
\[
\ln\!\left(n\binom{n}{h}\right) < \frac{2^h}{n},
\]
which implies
\[
n\binom{n}{h}e^{-2^h/n}<1.
\]
Therefore the bad-event probability in (3) is strictly less than $1$, so there exists a map
$f:\mathcal P(X)\to X$ for which no pair $(x,Y)$ is bad.

Thus every $h$-subset $Y$ satisfies $\{f(A):A\subseteq Y\}=X$, and therefore
\[
H(n)\le h=\left\lceil \log_2 n + 2\log_2\log_2 n + 2\right\rceil
\]
for all sufficiently large $n$. \hfill $\square$

\paragraph{Corollary.}
\[
H(n)\le \log_2 n + (2+o(1))\log_2\log_2 n.
\]

This is strictly stronger than the Round-1 quoted upper bound
\[
H(n)<\log_2 n + (3+o(1))\log_2\log_2 n.
\]

\subsection*{6) ADVERSARIAL VERIFICATION}

\begin{itemize}
\item \emph{Do larger sets $Y$ with $|Y|>h$ also work?}
Yes. If every $h$-subset sees all colours, then every larger set does too, since it contains an $h$-subset.

\item \emph{Did the random proof use independence correctly?}
Yes. For a fixed pair $(x,Y)$, the values $\{f(A):A\subseteq Y\}$ are independent by construction,
so the missing-colour probability is exactly $(1-\frac1n)^{2^h}$.

\item \emph{Is the barrier theorem too strong?}
No. It is intentionally limited: it rules out only the class of arguments obtained from nonnegative
linear combinations of the elementary inequalities (1) with coefficients independent of the colour.
It does \emph{not} rule out more sophisticated multi-colour or higher-order arguments.
Indeed, Alon's $H(2^k)\ge k+1$ argument from Round-1 already lies outside this purely first-moment framework.

\item \emph{Does the new upper bound contradict the conjectured lower bound?}
No. The original problem asks whether $H(n)-\log_2 n\to\infty$.
The theorem above merely shows that the excess above $\log_2 n$ is at most $2\log_2\log_2 n+O(1)$.
The conjectured divergence could still hold inside that window.
\end{itemize}

\subsection*{7) FINAL}

\textbf{UNRESOLVED (BUT STRICTLY ADVANCED).}

Round-2 gives two genuine advances beyond Round-1:

\begin{enumerate}
\item a new unconditional upper bound
\[
H(n)\le \log_2 n + 2\log_2\log_2 n + O(1),
\]
which improves the Round-1 quoted $(3+o(1))\log_2\log_2 n$ upper bound;

\item a rigorous barrier theorem showing that a broad first-moment / averaging class of lower-bound
arguments can never beat the trivial bound $n\le 2^h$.
\end{enumerate}

So the unresolved part of Problem 624 is now sharper:
to prove $H(n)-\log_2 n\to\infty$, one needs a genuinely structural lower-bound mechanism.

\subsection*{8) COMPLETION ESTIMATE (MANDATORY)}

COMPLETION: 45\%

\subsection*{9) REFERENCES}

\begin{itemize}
\item T. F. Bloom, \emph{Erd\H{o}s Problem \#624}, \texttt{erdosproblems.com/624}, accessed 2026-03-07.
\end{itemize}

\subsection*{10) OUTPUT FORMAT}

This section is written in drop-in \LaTeX\ form and starts directly from \verb|\section{}|, with no preamble.

\bigskip

